\documentclass[conference]{IEEEtran}
\usepackage{cite}
\usepackage{amsmath,amssymb,amsfonts}
\usepackage{graphicx}
\usepackage{textcomp}
\usepackage{xcolor}

\begin{document}

\title{Real-Time Collaborative Editing Without a Central Lock}

\author{\IEEEauthorblockN{First Author}
\IEEEauthorblockA{\textit{Dept. of Computer Science and Engineering} \\
\textit{University of Chittagong}\\
Chattogram, Bangladesh \\
first.author@example.com}
\and
\IEEEauthorblockN{Second Author}
\IEEEauthorblockA{\textit{Dept. of Computer Science and Engineering} \\
\textit{University of Chittagong}\\
Chattogram, Bangladesh \\
second.author@example.com}
}

\maketitle

\begin{abstract}
We describe a collaborative editor whose document is a sequence CRDT. Concurrent edits from
any number of users, including users working offline, merge deterministically without
coordination, and concurrently typed words never interleave.
\end{abstract}

\begin{IEEEkeywords}
CRDT, collaborative editing, distributed systems, eventual consistency
\end{IEEEkeywords}

\section{Introduction}
Operational transformation requires a central server to order operations~\cite{ellis1989}.
CRDTs instead guarantee that replicas which have seen the same set of operations hold the same
state~\cite{shapiro2011}.

\section{Method}
Each character carries a unique identifier $(a, s)$ where $a$ is the author and $s$ a counter.
Its position is determined by its left and right origins:
\begin{equation}
  \mathrm{pos}(c) = f\bigl(\mathrm{origin}_L(c), \mathrm{origin}_R(c), \mathrm{id}(c)\bigr).
\end{equation}

\section{Evaluation}
Replaying a 100{,}000-keystroke trace takes well under a second on a laptop.

\section{Conclusion}
CRDTs make offline-first collaborative writing practical.

\begin{thebibliography}{00}
\bibitem{ellis1989} C. A. Ellis and S. J. Gibbs, ``Concurrency control in groupware systems,'' in \textit{Proc. ACM SIGMOD}, 1989, pp. 399--407.
\bibitem{shapiro2011} M. Shapiro, N. Pregui\c{c}a, C. Baquero, and M. Zawirski, ``Conflict-free replicated data types,'' in \textit{Proc. SSS}, 2011, pp. 386--400.
\end{thebibliography}

\end{document}
